% ECONOMICS_PROBLEM: {"id": "E62-583", "title": "Fiscal credibility and expectations", "jel_code": "E62", "source_id": "OP-0260"}
% !TeX program = xelatex
\documentclass[11pt,a4paper]{article}
\usepackage[margin=23mm]{geometry}
\usepackage{fontspec}
\setmainfont{Latin Modern Roman}
\usepackage{amsmath,amssymb,mathtools}
\usepackage{xcolor,enumitem,booktabs,longtable,array,xurl,hyperref}
\definecolor{muted}{HTML}{53636C}
\hypersetup{colorlinks=true,urlcolor=blue,linkcolor=blue}
\setlength{\parindent}{0pt}
\setlength{\parskip}{5pt}
\newcommand{\R}{\mathbb R}
\newcommand{\E}{\mathbb E}
\newcommand{\N}{\mathbb N}
\newcommand{\ind}{\mathbf 1}
\newcommand{\Var}{\operatorname{Var}}
\newcommand{\Cov}{\operatorname{Cov}}
\newcommand{\supp}{\operatorname{supp}}
\DeclareMathOperator*{\argmin}{arg\,min}
\DeclareMathOperator*{\argmax}{arg\,max}
\newcommand{\entrymeta}[3]{{\sffamily\small\color{muted}\textbf{\texttt{#1}}\quad|\quad #2\par #3\par}}
\newcommand{\Problemref}[1]{\texttt{#1}}
\newcommand{\JELref}[2]{\texttt{#2} (JEL \texttt{#1})}
\begin{document}
\section*{Fiscal credibility and expectations}
\textbf{Problem ID:} \texttt{E62-583}\quad \textbf{JEL:} \texttt{E62}\par
% BEGIN ECONOMICS STATEMENT
\label{problem:OP-0260}
{\sffamily\small\color{muted}Original subject group: Fiscal policy and sovereign debt\par}
\label{definitions:problem:30007541}
\textbf{Audit classification:} Published research agenda. A fiscal expectations evidence gap is explicit, but the source supplies findings; the experiment is editorial. Verification concerns the cited version, not first priority or proof that this exact editorial specification remains unsolved on October 8, 2026.
\subsection*{Definitions and precise research target}
\paragraph{Editorial mathematical specification.} Fix a fiscal authority’s announced future tax or spending path \(A\), holding immediate implemented policy fixed. Randomly or quasi-randomly vary credible communication \(Z_i\) among households, firms, or investors exposed to the same path. Let \(B_i(Z_i)\) record beliefs about implementation, future disposable income, and debt sustainability; let \(X_i(Z_i)\) be measured spending, investment, or required financing yield. Define the announcement intention-to-treat effect
\[
\tau_X=E[X_i(1)-X_i(0)],
\]
and, for a predetermined credibility measure \(c_i\), its conditional counterpart \(E[X_i(1)-X_i(0)\mid c_i=c]\).
Identify these effects separately from concurrent fiscal actions and news about economic fundamentals. For a belief-mediated effect, state additional exclusion and sequential-ignorability assumptions, or provide bounds when communication directly changes attention or uncertainty. Measure persistence and compare promised with implemented policy so that credibility is not retrospectively defined by favorable outcomes. To infer sovereign financing effects, specify market spillovers and whether treatment changes the information held by marginal investors. Assess whether estimates transport across institutions and fiscal regimes. The target is causal expectation and credibility effects under a declared announcement design, not the claim that fiscal communication has no known effects or that one observed spread response identifies credibility.

Define credibility before treatment, for example \(c_i\) as the reported probability that a fixed announced fiscal path will be implemented. It is a subjective implementation belief, not a favorable realized fiscal outcome. For random assignment, consistency, and no spillovers, the total intention-to-treat effect is a difference in observed means; heterogeneous effects require overlap within predeclared credibility strata. A belief-mediated causal effect additionally needs a defined intervention on the mediator or a potential-outcome structural model and appropriate cross-world or sequential conditions; assignment alone supplies no such identification. For bonds, use yield changes or a pricing functional with maturity and risk compensation specified. Marketwide investor information creates interference, so household experiments cannot automatically identify sovereign-price responses.
\subsection*{Variants and assumption changes}
The following are \emph{editorial variants}, not additional published open conjectures.
\begin{enumerate}
\item \emph{Private announcement delivery.} Randomize exposure to a common public fiscal plan among households while holding actual taxes and transfers fixed. Identify consumption and implementation-belief effects, measuring persistence and attrition.
\item \emph{Market information shock.} Use an independently timed audited fiscal disclosure affecting marginal investors. Estimate sovereign-yield responses with concurrent monetary and macro news excluded, and specify equilibrium information spillovers; no individual-level no-interference assumption is imposed.
\end{enumerate}
\subsection*{References and source audit}
\begin{itemize}
\item Nicolas End; Gee Hee Hong. \emph{Trust What You Hear: Policy Communication, Expectations, and Fiscal Credibility}. 2022. Locator: Section1, printed p.10, fiscal expectations comparison; Sections3--5 provide existing empirical results. Primary text checked. \url{https://www.imf.org/-/media/files/publications/wp/2022/english/wpiea2022036-print-pdf.pdf}.
\end{itemize}
A fiscal expectations evidence gap is explicit, but the source supplies findings; the experiment is editorial.
\begin{enumerate}\raggedright
\item\hypertarget{definitions:source:30007541:1}{} Nicolas End; Gee Hee Hong. 2022. \emph{Trust What You Hear: Policy Communication, Expectations, and Fiscal Credibility}. Section1, printed p.10, fiscal expectations comparison; Sections3--5 provide existing empirical results.
\url{https://www.imf.org/-/media/files/publications/wp/2022/english/wpiea2022036-print-pdf.pdf}.
DOI: \url{https://doi.org/10.5089/9798400200748.001}.
\end{enumerate}

\subsection*{Common conventions: Source mathematical conventions}

These conventions apply to each entry unless explicitly replaced.

\textbf{Models and identification.} A primitive model \(m\) specifies all probability laws, feasible choices, information, equilibrium and measurement rules used by its target. The admissible class \(\mathcal M\) is fixed independently of outcomes. If \(P_O(m)\) is its observable probability law and \(\tau(m)\) the target, its identified set at observable law \(P\) is
\[
 \mathcal I_\tau(P)=\{\tau(m):m\in\mathcal M,\ P_O(m)=P\}.
\]
Point identification means that this set is a singleton. An empty set signals incompatibility of data and assumptions. Coordinate bounds need not describe the joint set. Bounds are sharp only if every retained target is attainable or arbitrarily approachable by compatible models. Finite bounds require explicit restrictions on unobserved quantities; unrestricted confounding, hidden wealth, extrapolation or arbitrary equilibrium selection can yield uninformative sets.

\textbf{Causal interventions.} A potential outcome \(Y_i(p)\) is the outcome under a complete policy regime \(p\), including timing, financing, information and market spillovers. The target population and units are fixed before treatment. With interference, use the complete assignment vector or a justified exposure mapping; individual no-interference notation alone cannot identify an economy-wide intervention. A difference of mean potential outcomes does not require the cross-world joint distribution. Conditioning on post-treatment migration, survival, enrollment or employment changes the estimand and needs separate justification.

\textbf{Statistical statements.} An estimator, confidence set, forecast or test specifies a sampling law, model class and loss/coverage criterion. Uniform \((1-\alpha)\)-coverage means \(\inf_{m\in\mathcal M}\Pr_m\{\tau(m)\in C_n\}\geq1-\alpha\), or an explicitly asymptotic version. The whole parameter space trivially has coverage, so informativeness requires a stated width, risk or power objective. A forecasting improvement is relative to a named benchmark, information set and evaluation population.

\textbf{Equilibrium and optimization.} An equilibrium comprises feasible plans, optimality, beliefs where needed, budget constraints and all market/resource clearing. The primitives or a stated benchmark define each correspondence \(\mathcal E(p;m)\). Nonemptiness is assumed or proved, never inferred just from compact policy sets. A selected-equilibrium target uses a declared selection rule; a robust target quantifies over all permitted equilibria. Use suprema or infima unless attainment is established. An \(\varepsilon\)-optimal policy is feasible and within \(\varepsilon\) of the finite optimum value. Report an empty feasible set separately.

\textbf{Welfare and accounting.} Normative utility, interpersonal normalization, social weights, numeraire, discounting and the use of public receipts are primitives. Behavioral utility need not coincide with the welfare criterion. Household welfare includes ownership income and rents once; adding profits again to compensating variation generally double-counts them. A monetary equivalent is defined by an explicit transfer and price convention, with monotonicity and range conditions. Average effects, mechanism contributions, transfers and welfare losses are different targets. Infinite sums require convergence and dynamic feasible plans require terminal or no-Ponzi conditions.

\textbf{Source versions.} A working-paper date, revision date and journal publication year are distinguished. A DOI for a journal version is not automatically the DOI of an earlier manuscript. Locators refer to the stated version and printed pages where specified. A metadata-only check does not verify a claim about a full-text conclusion. Every entry's source subsection narrows the provenance claim accordingly.
% END ECONOMICS STATEMENT
\end{document}
